../cictikz/src/cictikz/data/tex/cictikz_preamble.tex

% The same cross section with the parasitic bipolars drawn in
% (l02_esd figure 16): the lateral npn (NMOS source / substrate /
% n-well) and the vertical pnp (PMOS source / n-well / substrate),
% with the substrate and well resistances that close the thyristor
% loop.

\definecolor{echarge}{RGB}{85,155,235}
\definecolor{hcharge}{RGB}{160,25,35}
\definecolor{oxteal}{RGB}{140,195,195}

\begin{circuitikz}[american, thick, transform shape]
  ../cictikz/src/cictikz/data/tex/cictikz_lib.tex
  %%%%%%%%%%%%%%%%%%%%%%%%%%%%%%%%%%%%%%%%%%%%%%%%%%%%%%%%%%%%%%%%%%%%%%
%% Shared frame for the ESD build-up sequence in l02_esd.
%%
%% Four figures in the chapter are the same picture with more added each
%% time: a chip outline with three rails brought out to pads, then one
%% diode, then two, then the whole network with the grounded gate NMOS.
%% They only teach anything if they are literally the same drawing, so
%% the frame lives here and each figure adds its own devices.
%%
%% Coordinates: pads at x = 0, rails run to \esdRight. VDD at \esdVdd,
%% the signal pin at \esdPin, VSS at zero.
%%%%%%%%%%%%%%%%%%%%%%%%%%%%%%%%%%%%%%%%%%%%%%%%%%%%%%%%%%%%%%%%%%%%%%

\newcommand{\esdVdd}{5.0}
\newcommand{\esdPin}{2.6}
\newcommand{\esdRight}{8.8}

% A bond pad: the crossed square the chapter's originals use.
\newcommand{\esdPad}[1]{%
  \draw (-0.17,{#1-0.17}) rectangle (0.17,{#1+0.17});
  \draw (-0.17,{#1-0.17}) -- (0.17,{#1+0.17});
  \draw (-0.17,{#1+0.17}) -- (0.17,{#1-0.17});
}

% The chip outline, three rails, their pads and their names. #1 is the
% name of the middle rail, because the overview figure calls it PIN and
% the protection figures treat it as one signal pin among many.
\newcommand{\esdframe}{%
  \draw[dotted, gray] (-0.8,-1.3) rectangle ({\esdRight+0.3},{\esdVdd+1.3});
  \foreach \y/\name/\num in {\esdVdd/VDD/1, \esdPin/PIN/2, 0/VSS/0} {
    \draw (0,\y) -- (\esdRight,\y);
    \esdPad{\y}
    \node[armygreen, anchor=south west, scale=0.85] at (0.05,{\y+0.2}) {\name};
    %- above the rail, not left of it: the zap wiring comes in on the
    %- left and the numbers were landing on top of it
    \node[red, anchor=south east, scale=0.9] at (-0.25,{\y+0.18}) {\num};
  }
}

% A vertical diode conducting upward, from (#1,#2) to (#1,#3), with its
% body centred at #4 and named #5. Every protection device in these
% figures is one of these; drawing them through one macro is what makes
% the reverse-biased-in-normal-operation argument in the text hold
% across all four pictures.
%
% The centre is explicit because a diode spanning VSS to VDD would
% otherwise put its symbol exactly on the pin rail it is meant to cross
% without connecting.
\newcommand{\esdDiode}[5]{%
  \draw (#1,#2) -- (#1,{#4-0.45});
  \draw (#1,{#4-0.45}) to [D] (#1,{#4+0.45});
  \draw (#1,{#4+0.45}) -- (#1,#3);
  \fill (#1,#2) circle (0.075);
  \fill (#1,#3) circle (0.075);
  \node[red, anchor=west, scale=0.8] at ({#1+0.32},#4) {#5};
}

% A grounded gate NMOS at x=#1, drain on the rail at #2, source on the
% rail at #3, body centred at #4, node id #5 (letters only - it names a
% TikZ node) and caption #6. The centre is explicit for the same reason
% the diode's is: left to the midpoint, a VDD to VSS device lands on the
% pin rail it is meant to cross without touching. The gate goes to the source through a resistor,
% which is the whole trick: the resistor does nothing at DC, so the
% device is off whenever the chip is running, but during a zap the
% substrate current it develops lifts the gate and turns the parasitic
% bipolar on.
\newcommand{\esdggn}[6]{%
  \node[nmos] (esdm#5) at ({#1+0.35},#4) {};
  \draw (esdm#5.D) -- ({#1+0.35},#2);
  \draw (esdm#5.S) -- ({#1+0.35},#3);
  \fill ({#1+0.35},#2) circle (0.075);
  \fill ({#1+0.35},#3) circle (0.075);
  \draw (esdm#5.G) -- ({#1-0.35},#4);
  \draw ({#1-0.35},#4) to [R] ({#1-0.35},{#3+0.5});
  \draw ({#1-0.35},{#3+0.5}) -- ({#1-0.35},#3) -- ({#1+0.35},#3);
  %- beside the device, not above it: above puts it on the rail the
  %- device crosses without connecting
  \node[red, anchor=west, scale=0.8, align=left]
    at ({#1+0.95},#4) {#6};
}

% The zap itself: current forced in at one pad, taken out at another.
\newcommand{\esdZapIn}[1]{%
  \draw (-2.7,#1) to [I] (-1.3,#1);
  \draw (-1.3,#1) -- (0,#1);
}
\newcommand{\esdZapOut}[1]{%
  \draw (0,#1) -- (-1.3,#1);
  \draw (-1.3,#1) \vground;
}

  \scrframe
  \scrarrows

  \begin{scope}[color=black!65, thick]
    % lateral npn: emitter at the NMOS source, base in the substrate,
    % collector the n-well
    \draw (3.55,-2.0) -- (3.55,-2.9);       % base bar
    \draw (3.0,-0.8) -- (3.0,-2.1) -- (3.55,-2.35);
    \draw[->] (3.55,-2.35) -- (3.05,-2.12); % emitter arrow (out of base)
    \draw (3.55,-2.6) -- (4.6,-3.1) -- (6.4,-3.1);
    % base resistance to the substrate contact
    \draw (3.55,-2.45) -- (2.6,-2.45);
    \draw (2.6,-2.45) -- (2.5,-2.25) -- (2.3,-2.65) -- (2.1,-2.25)
      -- (1.9,-2.65) -- (1.8,-2.45) -- (1.05,-2.45);
    \draw (1.05,-2.45) -- (1.05,-0.8);
    % vertical pnp: emitter at the PMOS source, base the n-well,
    % collector the substrate
    \draw (9.9,-2.2) -- (9.9,-3.1);         % base bar
    \draw (10.25,-0.8) -- (10.25,-2.2);
    \draw[->] (10.25,-2.2) -- (9.92,-2.42); % emitter arrow (into base)
    \draw (9.9,-2.9) -- (10.6,-3.5) -- (10.6,-4.4);
    % base resistance to the well tap
    \draw (9.9,-2.65) -- (10.9,-2.65);
    \draw (10.9,-2.65) -- (11.0,-2.45) -- (11.2,-2.85) -- (11.4,-2.45)
      -- (11.6,-2.85) -- (11.7,-2.65) -- (12.4,-2.65);
    \draw (12.4,-2.65) -- (12.4,-1.3);
  \end{scope}

\end{circuitikz}
\end{document}
